% Redrawn 2026-10-15; main labels 8 pt, notes 7.5 pt; actual fonts embedded in PDF.
% export-reference-width-mm: 112
% RL^2 protocol redrawn from Duan et al. (2016), arXiv:1611.02779v2.
% Final width 112 mm; theta is fixed at deployment, memory is not.
\begin{tikzpicture}[x=\linewidth,y=1mm,
 font=\figurefont\fontsize{8}{10}\selectfont,
 box/.style={rounded corners=1.8pt,draw=FigureStroke,fill=FigureBlueFill,line width=.6pt,
   minimum height=12mm,align=center,inner sep=1.5mm},
 flow/.style={-{Latex[length=2mm]},draw=FigureStroke,line width=.75pt}]
 \draw[draw=FigureStroke,dashed,rounded corners=1.8pt,line width=.6pt,fill=white]
   (.28,-24) rectangle (.98,23);
 \node[anchor=north] at (.63,21) {部署参数$\bm\theta$固定};
 \node[box,text width=23mm] (input) at (.115,0)
   {当前$o_t$\\前一步$a_{t-1}$\\$r_{t-1},\delta_{t-1}$};
 \node[box,text width=23mm,fill=FigureGreenFill] (memory) at (.49,0)
   {可变循环记忆\\$\bm m_t$};
 \node[box,text width=17mm] (action) at (.84,0)
   {策略输出\\$a_t$};
 \draw[flow] (input) -- (memory);
 \draw[flow] (memory) -- (action);
 \draw[flow] (memory.south) -- (.49,-16) -- (.34,-16)
   |- (memory.west);
 \node[anchor=west] at (.53,-16) {供下一步使用};
 \node[anchor=north,align=center] at (.50,-30)
   {回合结束：重置环境，保留同任务记忆\\
    任务切换：重置记忆，开始独立适应试验};
\end{tikzpicture}
